\section{The Sundays and Days}

Each Sunday below is given with its station, its opening words and its readings, the prayer or
text that leads the week, and what the Fathers saw in it. The Ember Days have a section of their
own after the Third Sunday, and the last days before Christmas close the part. Each section ends
with counsel for the week that follows.

\subsection{The First Sunday: \latin{Ad te levavi}}\label{sec:first-sunday}

In 2026 the First Sunday of Advent falls on 29 November. The season begins the evening before, at
First Vespers. The chapter of those Vespers is from Romans 13, the Epistle the Mass will read the
next morning, and begins \latin{Hora est} (gloss: it is the hour); the antiphon at the
Magnificat begins \latin{Ecce nomen Domini} (gloss: behold, the name of the Lord).

The station is at St Mary Major. The Mass opens \latin{Ad te levavi} (gloss: to thee I have
lifted up), and reads St Paul to the Romans (13:11--14) and the Gospel of the signs before the
end (Luke 21:25--33). Its Collect is the first of the season's four prayers that begin
\latin{Excita}:

\begin{studyframe}
\latin{Excita, quaesumus, Domine, potentiam tuam, et veni: ut ab imminentibus peccatorum
nostrorum periculis, te mereamur protegente eripi, te liberante salvari: Qui vivis et regnas cum
Deo Patre in unitate Spiritus Sancti, Deus: per omnia saecula saeculorum. R. Amen.}
\medskip

Stir up, O Lord, we pray Thee, Thy strength, and come among us, that whereas through our sins
and wickedness we do justly apprehend Thy wrathful judgments hanging over us, Thy bountiful grace
and mercy may speedily help and deliver us; Who livest and reignest with God the Father, in the
unity of the Holy Ghost, one God, world without end.\par
{\smallskip\noindent\small Latin: Triptych's editorial text of the 1962 proper, the wording of the Pustet
Missal of 1862 read against the 1962 typical edition (p.~1). English: tr. Bute, \textit{The Roman
Breviary} (1908), Winter, p.~218.\par}
\end{studyframe}

The Latin asks that from the dangers our sins bring upon us we may be \latin{eripi} and
\latin{salvari} (gloss: rescued, saved), \latin{te protegente} and \latin{te liberante} (gloss:
with thee protecting, with thee setting free). The Church asks three more times before Christmas
in prayers that begin with the same word. The Collect of the Second Sunday begins \latin{Excita,
Domine, corda nostra} (gloss: stir up our hearts, O Lord); the Collect of the Ember Friday begins
again with \latin{Excita}; and the Collect of the Fourth Sunday returns to the words of the
first and adds \latin{et magna nobis virtute succurre} (gloss: and with great power help us).

\subsubsection{The hour to rise from sleep}

The Epistle opens with the season's first summons:

\begin{quote}
And that, knowing the season, that it is now the hour for us to rise from sleep. For now our
salvation is nearer than when we believed.\par
{\smallskip\noindent\small ---Romans 13:11, Douay-Rheims (Challoner).\par}
\end{quote}

\noindent Its last two verses are the ones Augustine was reading when he was converted:

\begin{quote}
So quickly I returned to the place where Alypius was sitting; for there had I put down the
volume of the apostles, when I rose thence. I grasped, opened, and in silence read that paragraph
on which my eyes first fell,---``Not in rioting and drunkenness, not in chambering and wantonness,
not in strife and envying; but put ye on the Lord Jesus Christ, and make not provision for the
flesh, to fulfil the lusts thereof.'' No further would I read, nor did I need; for instantly, as
the sentence ended,---by a light, as it were, of security infused into my heart,---all the gloom of
doubt vanished away.\par
{\smallskip\noindent\small ---Augustine, \textit{Confessions} VIII.12.29, tr. J.~G. Pilkington,
NPNF, 1st ser., vol.~1 (Buffalo, 1886).\par}
\end{quote}

The Church opens Advent with the verses that ended Augustine's doubt. The vices Paul names before
the command, rioting and drunkenness, chambering and wantonness, strife and envying, are concrete
enough to examine on the first evening of the season.

\subsubsection{Lift up your heads}

The Gospel is Christ's own description of the end: signs in the sun, moon and stars, distress of
nations, men withering away for fear, and then the Son of Man coming in a cloud with great power
and majesty. The 1962 Office reads at Matins the opening of St Gregory the Great's homily on it:

\begin{quote}
Our Lord and Saviour wisheth to find us ready at His second coming. Therefore He telleth us what
will be the evils of the world as it groweth old, that He may wean our hearts from worldly
affections. Here we read what great convulsions will go before the end, that, if we will not fear
God in our prosperity, we may at least be scourged into fearing His judgment when it is at
hand.\par
{\smallskip\noindent\small ---Gregory the Great, \textit{Homiliae in Evangelia} 1.1, tr. Bute,
\textit{The Roman Breviary} (1908), Winter, p.~216.\par}
\end{quote}

Later in the homily Gregory comes to the verse that follows the signs:

\begin{quote}
But when these things begin to come to pass, look up and lift up your heads, because your
redemption is at hand.\par
{\smallskip\noindent\small ---Luke 21:28, Douay-Rheims (Challoner).\par}
\end{quote}

\noindent He reads it as Christ's consolation of the elect: when the world whose friends they are
not comes to its close, the redemption they seek is near. As the \textit{Catena Aurea} gives his
words:

\begin{quote}
For in holy Scripture the head is often put for the mind, for as the members are ruled by the
head, so are the thoughts regulated by the mind. To lift up our heads then, is to raise up our
minds to the joys of the heavenly country.\par
{\smallskip\noindent\small ---Gregory the Great, \textit{Homiliae in Evangelia} 1.3, as the
\textit{Catena Aurea} gives it on Luke 21:28, Oxford tr. (1865), pp.~688--689.\par}
\end{quote}

St Leo the Great, preaching on the December fast in a sermon the Office read on this Sunday
before 1960, took a verse from the same discourse of Christ, ``Take heed lest haply your hearts
be overcharged with surfeiting and drunkenness, and care of this life,'' and pressed it on his
own people: ``And assuredly, dearly beloved, we acknowledge that this precept applies more
especially to us, to whom undoubtedly the day denounced is near, even though hidden''
(\textit{Sermo} XIX.1, tr. Feltoe). For that coming, he goes on, every man must prepare
himself, lest it find him given over to gluttony or entangled in the cares of this life; and the
preparation he asks of them in that sermon is the December fast.

\medskip
\noindent\textit{Counsel for the week.} Say the Collect \latin{Excita} each morning this week, in
Latin or in Bute's English. Read Romans 13:11--14 on the Sunday evening and name before God one
work of darkness to put off before Christmas. Begin the Isaiah lessons of
Appendix~\ref{sec:daily} on the Monday; the first chapters are an indictment of a people who know
God and do not act on what they know, and they repay slow reading. Once a day, lift up your head
from whatever presses, in Gregory's sense: raise the mind, for a moment, to the joys of the
heavenly country. Those who keep the St Andrew prayer of Appendix~\ref{sec:prayers} begin it on
30 November, the feast of St Andrew, and say it until Christmas (Section~\ref{sec:prayer}).
